% !TeX root = ../sections/03_solutions.tex

\subsection{2_1}
\begin{solution}
	与えられた関数は
	\[
	f(x)=
	\begin{cases}
		0
		& \left(|x|\leq \dfrac{\pi}{3}\right),\\[6pt]
		\dfrac{\pi}{2}
		& \left(\dfrac{\pi}{3}<|x|\leq \dfrac{2\pi}{3}\right),\\[6pt]
		\pi
		& \left(\dfrac{2\pi}{3}<|x|\leq \pi\right)
	\end{cases}
	\]
	であり，周期 \(2\pi\) で拡張されている。
	
	\medskip
	
	まずグラフについて考える。
	
	この関数は \(|x|\) の値だけで決まるので偶関数である。
	したがって，グラフは \(y\)-軸対称である。
	
	区間ごとの値は，
	\[
	|x|\leq \frac{\pi}{3}
	\quad\text{では}\quad
	f(x)=0,
	\]
	\[
	\frac{\pi}{3}<|x|\leq \frac{2\pi}{3}
	\quad\text{では}\quad
	f(x)=\frac{\pi}{2},
	\]
	\[
	\frac{2\pi}{3}<|x|\leq \pi
	\quad\text{では}\quad
	f(x)=\pi
	\]
	である。
	
	したがって，\(-\pi\leq x\leq \pi\) では，
	中央が高さ \(0\)，その外側が高さ \(\pi/2\)，さらに外側が高さ \(\pi\) の階段状のグラフになる。
	
	\medskip
	
	次に，フーリエ係数を求める。
	
	周期 \(2\pi\) のフーリエ係数は
	\[
	a_0=\frac{1}{\pi}\int_{-\pi}^{\pi}f(x)\,dx,
	\]
	\[
	a_n=\frac{1}{\pi}\int_{-\pi}^{\pi}f(x)\cos nx\,dx,
	\]
	\[
	b_n=\frac{1}{\pi}\int_{-\pi}^{\pi}f(x)\sin nx\,dx
	\]
	である。
	
	この関数は偶関数なので，
	\[
	b_n=0
	\]
	である。
	
	また，
	\[
	a_0
	=
	\frac{2}{\pi}\int_0^\pi f(x)\,dx
	\]
	である。
	
	区間ごとに分けると，
	\[
	a_0
	=
	\frac{2}{\pi}
	\left(
	\int_0^{\pi/3}0\,dx
	+
	\int_{\pi/3}^{2\pi/3}\frac{\pi}{2}\,dx
	+
	\int_{2\pi/3}^{\pi}\pi\,dx
	\right).
	\]
	
	これを計算すると，
	\[
	a_0
	=
	\frac{2}{\pi}
	\left(
	\frac{\pi}{2}\cdot \frac{\pi}{3}
	+
	\pi\cdot \frac{\pi}{3}
	\right)
	=
	\frac{2}{\pi}
	\left(
	\frac{\pi^2}{6}+\frac{\pi^2}{3}
	\right)
	=
	\frac{2}{\pi}\cdot \frac{\pi^2}{2}
	=
	\pi.
	\]
	
	したがって，
	\[
	a_0=\pi
	\]
	である。
	次に \(a_n\) を求める。
	偶関数なので，
	\[
	a_n
	=
	\frac{2}{\pi}\int_0^\pi f(x)\cos nx\,dx
	\]
	である。
	
	区間ごとに分けると，
	\[
	a_n
	=
	\frac{2}{\pi}
	\left(
	\int_{\pi/3}^{2\pi/3}
	\frac{\pi}{2}\cos nx\,dx
	+
	\int_{2\pi/3}^{\pi}
	\pi\cos nx\,dx
	\right).
	\]
	
	ここで，
	\[
	\int \cos nx\,dx
	=
	\frac{\sin nx}{n}
	\]
	を用いる。
	
	したがって，
	\[
	a_n
	=
	\frac{2}{\pi}
	\left(
	\frac{\pi}{2}
	\left[
	\frac{\sin nx}{n}
	\right]_{\pi/3}^{2\pi/3}
	+
	\pi
	\left[
	\frac{\sin nx}{n}
	\right]_{2\pi/3}^{\pi}
	\right).
	\]
	
	よって，
	\[
	a_n
	=
	\frac{2}{\pi}
	\left(
	\frac{\pi}{2n}
	\left(
	\sin\frac{2n\pi}{3}
	-
	\sin\frac{n\pi}{3}
	\right)
	+
	\frac{\pi}{n}
	\left(
	\sin n\pi
	-
	\sin\frac{2n\pi}{3}
	\right)
	\right).
	\]
	
	ここで，
	\[
	\sin n\pi=0
	\]
	であるから，
	\[
	a_n
	=
	\frac{2}{\pi}
	\left(
	\frac{\pi}{2n}
	\left(
	\sin\frac{2n\pi}{3}
	-
	\sin\frac{n\pi}{3}
	\right)
	-
	\frac{\pi}{n}
	\sin\frac{2n\pi}{3}
	\right).
	\]
	
	整理すると，
	\[
	a_n
	=
	\frac{1}{n}
	\left(
	\sin\frac{2n\pi}{3}
	-
	\sin\frac{n\pi}{3}
	\right)
	-
	\frac{2}{n}
	\sin\frac{2n\pi}{3}.
	\]
	
	したがって，
	\[
	a_n
	=
	-\frac{1}{n}
	\left(
	\sin\frac{n\pi}{3}
	+
	\sin\frac{2n\pi}{3}
	\right).
	\]
	
	以上より，フーリエ係数は
	\[
	\boxed{
		a_0=\pi,
		\qquad
		a_n=
		-\frac{1}{n}
		\left(
		\sin\frac{n\pi}{3}
		+
		\sin\frac{2n\pi}{3}
		\right),
		\qquad
		b_n=0
	}
	\]
	である。
	
	したがって，第 \(N\) 部分和は
	\[
	S_N(x)
	=
	\frac{\pi}{2}
	+
	\sum_{n=1}^{N}
	a_n\cos nx
	\]
	であり，
	\[
	a_n=
	-\frac{1}{n}
	\left(
	\sin\frac{n\pi}{3}
	+
	\sin\frac{2n\pi}{3}
	\right)
	\]
	を代入すればよい。
	\end{solution}
	\begin{solution}
		gnuplot では，まず任意の \(x\) を基本区間 \([-\pi,\pi)\) に戻す関数を定義する。
		
		\begin{verbatim}
			pi = 3.141592653589793
			
			p(x) = x - 2*pi*floor((x + pi)/(2*pi))
			
			f(x) = (abs(p(x)) <= pi/3) ? 0 : \
			(abs(p(x)) <= 2*pi/3) ? pi/2 : pi
			
			a(k) = -(sin(k*pi/3) + sin(2*k*pi/3))/k
			
			S(N,x) = pi/2 + sum [k=1:N] a(k)*cos(k*x)
			
			set yrange [-0.5:4]
			
			plot f(x) title "f(x)" linecolor 1, \
			S(7,x) title "S7(x)" linecolor 2, \
			S(31,x) title "S31(x)" linecolor 3
		\end{verbatim}
		
		これを \texttt{plot.txt} として保存する。
		
		ただし，問題文の指示により，
		\texttt{plot.txt} には \texttt{set terminal} や \texttt{set output} は含めない。
		
		画像ファイル \texttt{graph.png} を作成するには，
		gnuplot 上で次を実行する。
		
		\begin{verbatim}
			set terminal pngcairo size 1000,700
			set output "graph.png"
			load "plot.txt"
			set output
		\end{verbatim}
	\end{solution}